\bookchapter{Capítulo 47: Manual de consulta rápida e recursos}

\chapterintro{Este capítulo é um manual de consulta rápida que funciona por si só. Quando você não lembrar de algum princípio ou precisar tomar uma decisão rápida, vá direto para cá. Recomendo tratá-lo como um “cartão de bolso”: consulta dos cinco grandes princípios, árvore de decisão de ferramentas, perguntas frequentes e recursos de estudo — sem pular nenhuma página.}

\booksection{Cartão de consulta rápida dos modelos mentais}

\begin{center}
\begin{tabularx}{\linewidth}{|X|X|X|}
\hline
\textbf{Princípio} & \textbf{Em uma frase} & \textbf{Pontos práticos} \tabularnewline
\hline
Think Before Coding & Pense bem antes de pedir & Responda a três perguntas: o que fazer? como fica o resultado? quais são as restrições? \tabularnewline
Simplicity First & Resolva o problema com o mínimo de código & Faça um recurso por vez e corte a complexidade desnecessária \tabularnewline
Surgical Changes & Modifique com precisão, sem extrapolar & Defina claramente o escopo da alteração e recuse o "aproveitei e mudei mais um pouco" da IA \tabularnewline
Goal-Driven Execution & Diga "o que fazer", não "como fazer" & Dê o objetivo, não os passos, e deixe a IA escolher o caminho \tabularnewline
Iterate \& Verify & Ciclo de gerar, validar e corrigir & Validação em três camadas: funcionalidade, limites e arquitetura \tabularnewline
\hline
\end{tabularx}
\end{center}

\begin{center}\small\color{gray}Cartão de consulta rápida dos cinco grandes princípios de pensamento\end{center}

\booksection{Árvore de decisão para a escolha de ferramentas}

Quando você precisar escolher uma ferramenta de programação com IA, siga o caminho de decisão abaixo. Atenção: isto não é um ranking de “qual é a melhor”, e sim um quadro de julgamento para “qual é a mais adequada a você”.

O que você quer fazer?

Zero experiência, quer experimentar rápido escolha um Builder na nuvem (v0/Bolt/Replit)

Desenvolvimento do dia a dia, precisa de completamento de código escolha um IDE com IA embutida (Cursor/Trae)

Projetos grandes, precisa de operações em lote escolha um Agent no terminal (Claude Code)

Validar uma ideia com protótipo rápido escolha um Builder na nuvem ou o Trae SOLO

\begin{center}\small\color{gray}Figura 6: Árvore de decisão para escolher ferramentas de programação com IA\end{center}

\booksection{Perguntas frequentes (FAQ)}

P1: Não sei nada de programação. Posso programar com IA?

Pode. As ferramentas de programação com IA de 2026 já são fortes o bastante para você descrever a necessidade em linguagem natural e deixar a IA gerar o código. Mas dois conselhos: primeiro, comece por ferramentas Builder na nuvem (como v0 e Bolt), que não exigem configurar ambiente de desenvolvimento; segundo, aprenda os conceitos básicos enquanto usa — não é preciso dominar a sintaxe, mas você precisa conseguir “entender o sentido geral” do código que a IA escreve; só assim a validação do resultado é eficaz.

P2: Qual linguagem de programação devo aprender?

Depende do que você quer fazer. Front-end de sites — JavaScript/TypeScript; processamento de dados ou IA — Python; back-end corporativo — Java ou Go. Mas o mais importante é: na era da programação com IA, a importância da “linguagem” está caindo. A IA cuida dos detalhes de sintaxe; o que você precisa entender é a “estrutura dos programas” — o que é uma função, o que é um laço, o que são estruturas de dados. Esses conceitos atravessam as linguagens; uma vez compreendidos, trocar de linguagem é rápido.

P3: Código escrito por IA tem problema de direitos autorais?

Até agosto de 2026, ainda é uma questão jurídica controversa. Em linhas gerais, o código gerado por IA não é considerado com um titular de direitos autorais claro na maioria das jurisdições. Mas, se você for usar em projetos comerciais, recomenda-se: 1) conhecer os termos de propriedade intelectual nos termos de serviço da ferramenta de IA que você usa; 2) revisar e modificar manualmente o código essencial, de modo que ele receba o seu aporte criativo; 3) evitar que a IA copie diretamente código aberto que pode ter direitos autorais.

P4: A IA vai substituir os programadores?

Não em curto prazo. O que a IA substitui não é a profissão de “programador”, e sim o trabalho de “codificar linha por linha”. Assim como a calculadora não substituiu os matemáticos, mas os libertou do cálculo manual maçante. As ferramentas de programação com IA deslocam o programador de “escrever código” para “pensar nos problemas” — o que, na prática, eleva o teto de valor do programador. A competitividade do futuro não está em “quantas linhas de código você escreve”, e sim em “quão complexo é o problema que você resolve”.

P5: Como julgar se a entrega da IA está correta?

Pela validação em três camadas. Primeira: validação funcional — roda? as entradas e saídas estão certas? É o mais básico. Segunda: validação de limites — o que acontece com entrada vazia, entrada longa demais, entrada anormal? Isso revela problemas escondidos. Terceira: validação de arquitetura — a estrutura do código faz sentido? há valores fixados no código (hardcoded)? é fácil de ampliar? Aqui você precisa de certo julgamento técnico. Se você não consegue avaliar a terceira camada, cumpra pelo menos as duas primeiras — se roda e os limites se comportam bem, você já está à frente da maioria dos iniciantes.

\booksection{Recursos de estudo recomendados}

A seguir, os recursos de estudo recomendados até agosto de 2026. Atenção: links de internet podem mudar; se algum link deixar de funcionar, buscar pelo título normalmente leva à versão mais recente.

\begin{center}
\begin{tabularx}{\linewidth}{|X|X|X|}
\hline
\textbf{Tipo} & \textbf{Recurso} & \textbf{Por que recomendamos} \tabularnewline
\hline
Fonte original & CLAUDE.md do Andrej Karpathy & A fonte dos cinco grandes princípios; 70 linhas, todas preciosas \tabularnewline
Fonte original & Documentação oficial da Anthropic (Claude Code) & Guia de melhores práticas do agente de terminal \tabularnewline
Tutorial & Coleção de Skills organizada pela comunidade de Vibe Coding & Conjunto de técnicas de campo reunidas por praticantes da comunidade \tabularnewline
Tutorial & Documentação oficial do Cursor e canal no YouTube & O guia autoritativo de uso da IA embutida no IDE \tabularnewline
Comunidade & awesome-claude-code no GitHub & Melhores práticas e coleção de plugins do Claude Code \tabularnewline
Comunidade & Comunidade r/cursor no Reddit & Experiências compartilhadas e discussões de usuários do Cursor no mundo todo \tabularnewline
Básico & MDN Web Docs (developer.mozilla.org) & A documentação de referência mais autoritativa do desenvolvimento front-end \tabularnewline
Básico & Tutorial oficial de Python (docs.python.org) & O ponto de partida oficial para aprender Python \tabularnewline
\hline
\end{tabularx}
\end{center}

\begin{center}\small\color{gray}Lista de recursos de estudo recomendados\end{center}

\booksection{Palavras finais: ferramentas nas mãos, metodologia na cabeça}

Voltemos à crença central deste livro: as ferramentas iteram a cada três meses; a metodologia evolui a cada três anos. Hoje, em agosto de 2026, as ferramentas de programação com IA já são poderosas a ponto de impressionar — você cria um site com uma frase, gera um aplicativo full stack com um parágrafo de descrição e conclui uma refatoração complexa com um único comando. Mas quanto mais poderosas as ferramentas, mais se evidencia a importância da metodologia — porque o poder das ferramentas amplifica cada uma das suas decisões.

Se o seu pedido é claro, a IA implementa com eficiência para você; se o seu pedido é vago, a IA também produz a coisa errada com eficiência. Se você segue o “simplicidade em primeiro lugar”, a IA ajuda a validar as suas ideias depressa; se você quer abraçar o mundo, a IA também acumula dívida técnica depressa. Se você modifica com precisão, a IA otimiza o código como uma cirurgia; se a sua instrução é vaga, a IA também instala bugs a golpes largos.

É por isso que este livro dedicou a maior parte do espaço à metodologia, e não à operação das ferramentas. Não porque as ferramentas não importem, mas porque a metodologia é o “multiplicador” das ferramentas — quanto melhor a metodologia, maior o valor das ferramentas; com metodologia zero, por mais fortes que sejam as ferramentas, é zero vezes qualquer número.

Você não programa mais sozinho. Você é o comandante e a IA é seu exército; a qualidade do seu pensamento determina a força do seu exército. A partir de hoje, arme seu pensamento com metodologia e amplifique sua capacidade com ferramentas — esse é todo o segredo da programação com IA.

\booksection{Referências}

\begin{enumerate}
  \item Andrej Karpathy. (2025). Vibe Coding. Social media post introducing the concept of vibe-driven programming with AI.
  \item Andrej Karpathy. (2025). CLAUDE.md - Project guidelines for AI-assisted development. GitHub repository, 70+K stars.
  \item Anthropic. (2026). Claude Code: Terminal-native AI programming. Official documentation and best practices.
  \item Cursor Team. (2026). Cursor 3: The next generation of AI-first IDE. Release notes and documentation.
  \item OpenAI. (2026). Codex CLI: Bringing AI to the terminal. Technical documentation.
  \item Model Context Protocol (MCP) Specification. (2025-2026). Open protocol for AI-tool integration.
  \item Vibe Coding Community. (2026). Community-curated skills and best practices for AI-assisted programming.
  \item Spec-Driven Development (SDD) methodology. (2026). Emerging practices in specification-first AI programming.
\end{enumerate}
